\documentclass[11pt]{article}

\usepackage[T1]{fontenc}
\usepackage[utf8]{inputenc}
\usepackage{lmodern}
\usepackage{geometry}
\usepackage{hyperref}
\usepackage{enumitem}
\usepackage{booktabs}
\usepackage{array}
\usepackage{amsmath}
\usepackage{amssymb}
\usepackage{xcolor}

\geometry{margin=1in}
\hypersetup{colorlinks=true,linkcolor=blue,urlcolor=blue,citecolor=blue}
\setlist[itemize]{topsep=3pt,itemsep=2pt,parsep=0pt}
\setlist[enumerate]{topsep=3pt,itemsep=2pt,parsep=0pt}

% LaTeX assigns a numerical depth to each sectioning level:0 = Chapters, 1 = Sections, 2 = Subsections, 3 = Subsubsections
\setcounter{tocdepth}{2}

\begin{document}

\title{\Huge{Alphalens Documentation}}
\date{}

\maketitle

\tableofcontents

\newpage

\section{Overview}

Alphalens is organized as a pipeline of four core modules:

\begin{itemize}
  \item \texttt{alphalens.utils} prepares factor data and computes forward returns.
  \item \texttt{alphalens.performance} analyzes factor behavior and simulated portfolio performance.
  \item \texttt{alphalens.plotting} visualizes metrics and diagnostics.
  \item \texttt{alphalens.tears} assembles the results into tear sheets.
\end{itemize}

The usual workflow is:

\begin{enumerate}
  \item Build cleaned factor data with \texttt{alphalens.utils.get\_clean\_factor\_and\_forward\_returns()}.
  \item Compute diagnostics and portfolio statistics with \texttt{alphalens.performance}.
  \item Visualize outputs with \texttt{alphalens.plotting}.
  \item Generate presentation-ready reports with \texttt{alphalens.tears}.
\end{enumerate}

All modules typically operate on a MultiIndex \texttt{pandas} object indexed by \texttt{(date, asset)}.

\section{alphalens.utils}

Utilities for preparing factor data, computing forward returns, binning factors into quantiles/bins, and working with custom trading calendars.

This module is the main data-preparation layer used throughout Alphalens. Most workflows start with \texttt{get\_clean\_factor\_and\_forward\_returns()} and then feed the resulting MultiIndex \texttt{DataFrame} into the tear sheets and plotting functions.

\subsection{Data model used by this module}

Most functions operate on a \texttt{pandas} object indexed by:

\begin{itemize}
  \item level 0: \texttt{date} / timestamp
  \item level 1: \texttt{asset}
\end{itemize}

Forward returns are usually stored in columns named with \texttt{pd.Timedelta}-compatible strings such as \texttt{1D}, \texttt{5D}, \texttt{30m}, or \texttt{1D1h}.

\subsection{Exceptions}

\paragraph{\texttt{NonMatchingTimezoneError}}
Raised when factor and pricing indices have different timezones.

\paragraph{\texttt{MaxLossExceededError}}
Raised by \texttt{get\_clean\_factor()} when the share of dropped factor data exceeds \texttt{max\_loss}.

\subsection{Core workflow}

\subsubsection{\texttt{compute\_forward\_returns(factor, prices, periods=(1, 5, 10), filter\_zscore=None, cumulative\_returns=True)}}
Compute N-period forward returns for each asset in \texttt{factor}.

\paragraph{Input expectations}
\begin{itemize}
  \item \texttt{factor}: MultiIndex \texttt{Series} indexed by \texttt{(date, asset)}.
  \item \texttt{prices}: wide \texttt{DataFrame} with dates on the index and assets on the columns.
  \item The price history must extend far enough beyond the factor timestamps to cover the largest requested period.
\end{itemize}

\paragraph{Behavior}
\begin{itemize}
  \item Validates that factor and price timezones match.
  \item Infers a trading calendar when the factor index has no frequency.
  \item Computes forward returns for each requested period.
  \item Optionally filters out extreme values using \texttt{filter\_zscore}.
  \item Names the output columns using a \texttt{Timedelta}-style string.
\end{itemize}

\paragraph{Returns}
A MultiIndex \texttt{DataFrame} with the same \texttt{(date, asset)} index as the input factor.

\subsubsection{\texttt{get\_clean\_factor(factor, forward\_returns, groupby=None, binning\_by\_group=False, quantiles=5, bins=None, groupby\_labels=None, max\_loss=0.35, zero\_aware=False)}}
Align a factor series with already-computed forward returns and assign factor quantiles or bins.

Use this when you already have forward returns and do not need to recompute them.

\paragraph{Key behavior}
\begin{itemize}
  \item Joins factor values to forward returns.
  \item Optionally maps assets to groups.
  \item Computes factor quantiles or bins using \texttt{quantize\_factor()}.
  \item Drops rows with missing data.
  \item Tracks how much data was lost and raises \texttt{MaxLossExceededError} if the loss exceeds \texttt{max\_loss}.
\end{itemize}

\paragraph{Returns}
A cleaned MultiIndex \texttt{DataFrame} containing forward return columns, \texttt{factor}, optionally \texttt{group}, and \texttt{factor\_quantile}.

\subsubsection{\texttt{get\_clean\_factor\_and\_forward\_returns(factor, prices, groupby=None, binning\_by\_group=False, quantiles=5, bins=None, periods=(1, 5, 10), filter\_zscore=20, groupby\_labels=None, max\_loss=0.35, zero\_aware=False, cumulative\_returns=True)}}
One-stop helper that computes forward returns and then cleans and bins the factor data.

This is the primary entry point for most Alphalens analyses.

\paragraph{Equivalent steps}
\begin{enumerate}
  \item Call \texttt{compute\_forward\_returns()}
  \item Call \texttt{get\_clean\_factor()}
\end{enumerate}

\paragraph{Returns}
A cleaned MultiIndex \texttt{DataFrame} suitable for Alphalens tear sheets and plots.

\subsection{Quantization and bucketing}

\subsubsection{\texttt{quantize\_factor(factor\_data, quantiles=5, bins=None, by\_group=False, no\_raise=False, zero\_aware=False)}}
Assign factor values to quantile or value bins.

\paragraph{Behavior}
\begin{itemize}
  \item Exactly one of \texttt{quantiles} or \texttt{bins} must be provided.
  \item Supports group-wise binning with \texttt{by\_group=True}.
  \item Supports zero-aware binning, which splits positive and negative values separately.
  \item Can suppress binning failures with \texttt{no\_raise=True}, returning \texttt{NaN} for problematic rows.
\end{itemize}

\paragraph{Returns}
A \texttt{Series} named \texttt{factor\_quantile} indexed by \texttt{(date, asset)}.

\subsubsection{\texttt{non\_unique\_bin\_edges\_error(func)}}
Decorator that adds a more helpful error message when \texttt{pandas} cannot create unique bin edges.

This is used internally around \texttt{quantize\_factor()} to explain common binning failures, especially when many identical factor values span multiple quantiles.

\subsection{Return and performance helpers}

\subsubsection{\texttt{backshift\_returns\_series(series, N)}}
Shift a MultiIndex returns series backward by \texttt{N} observations in the first level.

This helper is useful for converting backward-looking returns into forward-looking returns.

\subsubsection{\texttt{demean\_forward\_returns(factor\_data, grouper=None)}}
Demean forward returns by date or by a custom grouper.

\paragraph{Behavior}
\begin{itemize}
  \item If \texttt{grouper} is not provided, demeaning happens per date.
  \item Only forward-return columns are adjusted.
  \item The result preserves the original shape and index.
\end{itemize}

\subsubsection{\texttt{rate\_of\_return(period\_ret, base\_period)}}
Convert returns observed over one period length into an equivalent rate for \texttt{base\_period}.

This is useful for normalizing returns across different horizons.

\subsubsection{\texttt{std\_conversion(period\_std, base\_period)}}
Convert standard deviation or standard error from one period length to another.

\subsection{Calendar and time utilities}

\subsubsection{\texttt{infer\_trading\_calendar(factor\_idx, prices\_idx)}}
Infer a trading calendar from factor and price datetimes.

\paragraph{Behavior}
\begin{itemize}
  \item Detects which weekdays are traded.
  \item Infers holidays by comparing the observed timestamps to a weekday-specific business-day calendar.
  \item Returns a \texttt{CustomBusinessDay} offset.
\end{itemize}

\subsubsection{\texttt{timedelta\_to\_string(timedelta)}}
Convert a \texttt{pd.Timedelta} into a compact string compatible with \texttt{pd.Timedelta(...)}.

Example output formats include \texttt{1D}, \texttt{3h15m}, and \texttt{1D1h}.

\subsubsection{\texttt{timedelta\_strings\_to\_integers(sequence)}}
Convert a sequence of timedelta strings into integer day counts.

Example: \texttt{['1D', '5D'] -> [1, 5]}.

\subsubsection{\texttt{add\_custom\_calendar\_timedelta(input, timedelta, freq)}}
Add a \texttt{Timedelta} to a timestamp or \texttt{DatetimeIndex} while respecting a custom trading calendar.

\paragraph{Accepted \texttt{freq} values}
\begin{itemize}
  \item \texttt{Day}
  \item \texttt{BusinessDay}
  \item \texttt{CustomBusinessDay}
\end{itemize}

\subsubsection{\texttt{make\_naive\_ts(t)}}
Return a timezone-naive timestamp.

\begin{itemize}
  \item If \texttt{t} is timezone-aware, it is converted to naive UTC-localized time.
  \item Otherwise it is localized to \texttt{None}.
\end{itemize}

\subsubsection{\texttt{diff\_custom\_calendar\_timedeltas(start, end, freq)}}
Compute the effective elapsed time between two timestamps under a custom calendar.

This is used when forward-return horizons must respect trading days, weekends, and holidays.

\paragraph{Accepted \texttt{freq} values}
\begin{itemize}
  \item any \texttt{pandas.tseries.offsets.BaseOffset}
  \item commonly \texttt{Day}, \texttt{BusinessDay}, or \texttt{CustomBusinessDay}
\end{itemize}

\subsection{Column detection and display helpers}

\subsubsection{\texttt{get\_forward\_returns\_columns(columns, require\_exact\_day\_multiple=False)}}
Identify which columns look like forward-return horizons.

\paragraph{Behavior}
\begin{itemize}
  \item Recognizes \texttt{Timedelta}-style labels such as \texttt{1D}, \texttt{5D}, \texttt{30m}, \texttt{1D1h}.
  \item When \texttt{require\_exact\_day\_multiple=True}, only exact day multiples are kept.
  \item Emits a warning if non-day-multiple columns are skipped in that mode.
\end{itemize}

\subsubsection{\texttt{print\_table(table, name=None, fmt=None)}}
Pretty-print a \texttt{Series} or \texttt{DataFrame}.

\paragraph{Behavior}
\begin{itemize}
  \item Uses rich display output when available.
  \item Falls back to standard formatted output.
  \item Temporarily changes \texttt{pandas} floating-point formatting when \texttt{fmt} is provided.
\end{itemize}

\subsection{Error handling helper}

\subsubsection{\texttt{rethrow(exception, additional\_message)}}
Re-raise an exception while preserving the original stack trace and appending extra context to the message.

This is used internally to make binning errors easier to diagnose.

\subsection{Practical usage}

Typical Alphalens usage looks like this:

\begin{enumerate}
  \item Prepare a factor series indexed by \texttt{(date, asset)}.
  \item Prepare a price \texttt{DataFrame} with assets in columns.
  \item Call \texttt{get\_clean\_factor\_and\_forward\_returns()}.
  \item Pass the result into tear sheets or plotting helpers.
\end{enumerate}

Example call: \texttt{get\_clean\_factor\_and\_forward\_returns(factor=..., prices=..., periods=(1, 5, 10))}.

\subsection{Notes}

\begin{itemize}
  \item Several functions assume data is already aligned on \texttt{(date, asset)}.
  \item Zero-aware binning is intended for factors centered around zero.
  \item \texttt{filter\_zscore} can introduce lookahead bias, so use it carefully.
  \item Forward-return horizon labels are derived from the price calendar and may reflect custom business-day offsets.
\end{itemize}

\section{alphalens.performance}

Performance-analysis utilities for factor research, including information coefficient calculations, factor-weighted portfolio simulation, quantile analysis, event-study helpers, and Pyfolio-ready outputs.

This module is the main analytics layer used after factor data has been cleaned and aligned with \texttt{alphalens.utils.get\_clean\_factor\_and\_forward\_returns()}.

\subsection{Data model used by this module}

Most functions expect a MultiIndex \texttt{pandas} object indexed by:

\begin{itemize}
  \item level 0: \texttt{date}
  \item level 1: \texttt{asset}
\end{itemize}

The input usually includes:

\begin{itemize}
  \item a \texttt{factor} column
  \item one or more forward-return columns such as \texttt{1D}, \texttt{5D}, or \texttt{10D}
  \item \texttt{factor\_quantile}
  \item optionally \texttt{group}
\end{itemize}

Forward-return column names are detected using \texttt{alphalens.utils.get\_forward\_returns\_columns()}.

\subsection{Core analysis functions}

\subsubsection{\texttt{factor\_information\_coefficient(factor\_data, group\_adjust=False, by\_group=False)}}
Compute the Spearman rank correlation between factor values and each forward-return horizon.

\paragraph{Behavior}
\begin{itemize}
  \item Calculates the Information Coefficient (IC) period by period.
  \item Optionally demeans forward returns by group before computing IC.
  \item Optionally computes IC separately for each group.
  \item Preserves the factor data frequency on the date index when not grouped by asset group.
\end{itemize}

\paragraph{Returns}
A \texttt{DataFrame} of IC values indexed by date, or by date and group when \texttt{by\_group=True}.

\subsubsection{\texttt{mean\_information\_coefficient(factor\_data, group\_adjust=False, by\_group=False, by\_time=None)}}
Compute the mean IC across the full sample or over a time window.

\paragraph{Behavior}
\begin{itemize}
  \item Delegates to \texttt{factor\_information\_coefficient()}.
  \item Can aggregate by time period using a pandas time rule such as monthly or weekly windows.
  \item Can aggregate by group when \texttt{by\_group=True}.
\end{itemize}

\paragraph{Returns}
A scalar-like \texttt{Series}/\texttt{DataFrame} of mean IC values depending on the requested grouping.

\subsection{Portfolio weighting and return simulation}

\subsubsection{\texttt{factor\_weights(factor\_data, demeaned=True, group\_adjust=False, equal\_weight=False)}}
Compute asset weights from factor values.

\paragraph{Behavior}
\begin{itemize}
  \item Factor values are normalized to gross leverage of 1.
  \item When \texttt{demeaned=True}, weights are constructed as a long-short portfolio.
  \item When \texttt{group\_adjust=True}, weights are computed in a group-neutral way.
  \item When \texttt{equal\_weight=True}, assets are equal-weighted instead of factor-weighted.
\end{itemize}

\paragraph{Returns}
A \texttt{Series} of weights indexed by \texttt{(date, asset)}.

\paragraph{Worked example}
Suppose one date contains three assets with these factor values:

\begin{center}
\begin{tabular}{lr}
\toprule
asset & factor \\
\midrule
A & 2.0 \\
B & 1.0 \\
C & -1.0 \\
\bottomrule
\end{tabular}
\end{center}

With the default \texttt{demeaned=True}, \texttt{group\_adjust=False}, and \texttt{equal\_weight=False}:

\begin{enumerate}
  \item Demean factor values by date.
  \begin{itemize}
    \item Mean factor = \((2.0 + 1.0 - 1.0) / 3 = 0.6667\)
    \item Demeaned values = \texttt{A: 1.3333}, \texttt{B: 0.3333}, \texttt{C: -1.6667}
  \end{itemize}
  \item Normalize by the sum of absolute values.
  \begin{itemize}
    \item Gross exposure = \(|1.3333| + |0.3333| + |-1.6667| = 3.3333\)
    \item Weights = \texttt{A: 0.40}, \texttt{B: 0.10}, \texttt{C: -0.50}
  \end{itemize}
\end{enumerate}

These weights represent a dollar-neutral long-short portfolio: long positions (\texttt{A: 0.40}, \texttt{B: 0.10}) total \(0.50\) and short positions (\texttt{C: -0.50}) total \(0.50\) in absolute value.

If \texttt{demeaned=False}, weights come directly from the raw (non-demeaned) factor values: the sum of absolute raw factors would be \(|2.0| + |1.0| + |-1.0| = 4.0\), yielding weights \texttt{A: 0.50}, \texttt{B: 0.25}, \texttt{C: -0.25}. If \texttt{group\_adjust=True}, weights are adjusted to be group-neutral. If \texttt{equal\_weight=True}, assets are equal-weighted within long and short buckets instead of using factor magnitudes.

\subsubsection{\texttt{factor\_returns(factor\_data, demeaned=True, group\_adjust=False, equal\_weight=False, by\_asset=False)}}
Compute period-wise returns for a factor-weighted portfolio.

In this module, a \textbf{factor return} is the return of a simulated portfolio whose asset weights are derived from the factor signal. For each date and each forward-return horizon, the return is computed as the sum of:

\[
\text{asset weight} \times \text{asset forward return}
\]

Equivalently:

\[
\text{factor return}[t, \text{period}] = \sum_i w_{t,i} \cdot r_{t,i,\text{period}}
\]

\paragraph{Behavior}
\begin{itemize}
  \item Uses \texttt{factor\_weights()} to build the portfolio.
  \item Multiplies weights by forward-return columns.
  \item If \texttt{by\_asset=True}, returns the weighted return contribution of each asset.
  \item Otherwise, aggregates to date-level portfolio returns.
\end{itemize}

\paragraph{Worked example}
Suppose one date contains three assets with these factor values and \texttt{1D} forward returns:

\begin{center}
\begin{tabular}{lrr}
\toprule
asset & factor & \texttt{1D} forward return \\
\midrule
A & 2.0 & 0.010 \\
B & 1.0 & 0.020 \\
C & -1.0 & -0.010 \\
\bottomrule
\end{tabular}
\end{center}

With the default \texttt{demeaned=True}, \texttt{group\_adjust=False}, and \texttt{equal\_weight=False}:

\begin{enumerate}
  \item Demean factor values by date.
  \begin{itemize}
    \item Mean factor = \((2.0 + 1.0 - 1.0) / 3 = 0.6667\)
    \item Demeaned values = \texttt{A: 1.3333}, \texttt{B: 0.3333}, \texttt{C: -1.6667}
  \end{itemize}
  \item Normalize by the sum of absolute values.
  \begin{itemize}
    \item Gross exposure = \(|1.3333| + |0.3333| + |-1.6667| = 3.3333\)
    \item Weights = \texttt{A: 0.40}, \texttt{B: 0.10}, \texttt{C: -0.50}
  \end{itemize}
  \item Multiply each weight by the forward return.
  \begin{itemize}
    \item \texttt{A: 0.40 \(\times\) 0.010 = 0.0040}
    \item \texttt{B: 0.10 \(\times\) 0.020 = 0.0020}
    \item \texttt{C: -0.50 \(\times\) -0.010 = 0.0050}
  \end{itemize}
  \item Sum across assets.
  \begin{itemize}
    \item factor return = 0.0040 + 0.0020 + 0.0050 = 0.0110
  \end{itemize}
\end{enumerate}

So the factor portfolio return for that date and period is \(0.011\), or 1.1\%.

If \texttt{demeaned=False}, the same formula is used, but the weights come directly from the raw factor values rather than demeaned values. If \texttt{group\_adjust=True}, weights are computed in a group-neutral way. If \texttt{equal\_weight=True}, the code uses equal-weighted long/short buckets instead of factor-magnitude weights.

\paragraph{Returns}
A \texttt{DataFrame} of period returns, or asset-level weighted returns when \texttt{by\_asset=True}.

\subsubsection{\texttt{factor\_alpha\_beta(factor\_data, returns=None, demeaned=True, group\_adjust=False, equal\_weight=False)}}
Estimate alpha and beta for a factor portfolio using OLS regression.

\paragraph{Behavior}
\begin{itemize}
  \item Uses factor portfolio returns as the dependent variable.
  \item Uses the universe mean return as the explanatory variable.
  \item If \texttt{returns} is not provided, it is computed using \texttt{factor\_returns()}.
  \item Annualizes alpha using a 252-trading-day convention.
\end{itemize}

\paragraph{Returns}
A \texttt{DataFrame} containing at least annualized alpha and beta by forward-return horizon.

\subsubsection{\texttt{cumulative\_returns(returns)}}
Convert simple returns into cumulative returns.

This is a thin wrapper around \texttt{empyrical.cum\_returns()} with a starting value of 1.

\subsubsection{\texttt{positions(weights, period, freq=None)}}
Build a time series of portfolio positions from factor weights.

\paragraph{Behavior}
\begin{itemize}
  \item Treats weights as active for a holding period defined by \texttt{period}.
  \item Uses the provided trading calendar frequency, or infers one from the weight index.
  \item Falls back to \texttt{BDay} and emits a warning if no frequency is available.
  \item Recomputes portfolio weights as positions roll forward through time.
\end{itemize}

\paragraph{Returns}
A \texttt{DataFrame} with timestamps on the index and assets on the columns.

\subsection{Quantile and bucket analysis}

\subsubsection{\texttt{mean\_return\_by\_quantile(factor\_data, by\_date=False, by\_group=False, demeaned=True, group\_adjust=False)}}
Compute mean forward returns and standard errors by factor quantile.

\paragraph{Behavior}
\begin{itemize}
  \item Can compute results by date or across the full sample.
  \item Can compute results by group.
  \item Can demean by the whole universe or within each group.
\end{itemize}

\paragraph{Returns}
A pair: \((\text{mean\_ret}, \text{std\_error\_ret})\).

\subsubsection{\texttt{compute\_mean\_returns\_spread(mean\_returns, upper\_quant, lower\_quant, std\_err=None)}}
Compute the difference in mean returns between two quantiles.

\paragraph{Behavior}
\begin{itemize}
  \item Subtracts lower-quantile mean returns from upper-quantile mean returns.
  \item Optionally propagates standard error for the spread.
\end{itemize}

\paragraph{Returns}
A pair: \((\text{mean\_return\_difference}, \text{joint\_std\_err})\).

\subsubsection{\texttt{quantile\_turnover(quantile\_factor, quantile, period=1)}}
Measure the proportion of names that leave a given quantile over time.

\paragraph{Behavior}
\begin{itemize}
  \item Compares membership in the selected quantile against the prior period.
  \item Preserves the date frequency of the input.
\end{itemize}

\paragraph{Worked example}
Suppose we evaluate \texttt{quantile=5} with \texttt{period=1}, and the names in quantile 5 are:

\begin{center}
\begin{tabular}{ll}
date & names in quantile 5 \\
2026-01-02 & \(\{A, B, C\}\) \\
2026-01-03 & \(\{B, C, D\}\) \\
2026-01-06 & \(\{C, D, E\}\) \\
\end{tabular}
\end{center}

For each date after the first, turnover is:

\[
\frac{\#\text{ new names vs previous date}}{\text{current quantile size}}
\]

\begin{itemize}
  \item On \texttt{2026-01-03}, new names vs \texttt{2026-01-02} are \(\{D\}\), so turnover is \(1 / 3 = 0.3333\).
  \item On \texttt{2026-01-06}, new names vs \texttt{2026-01-03} are \(\{E\}\), so turnover is \(1 / 3 = 0.3333\).
\end{itemize}

The first date has no prior comparison point, so no turnover value is reported for it.

\paragraph{Companion example (\texttt{period=2})}
Using the same memberships, each date is compared to the set two dates earlier.

\begin{itemize}
  \item On \texttt{2026-01-06}, compare \(\{C, D, E\}\) to \texttt{2026-01-02} (\(\{A, B, C\}\)).
  \item New names are \(\{D, E\}\), so turnover is \(2 / 3 = 0.6667\).
\end{itemize}

This highlights that larger \texttt{period} values measure non-adjacent membership change.
If intermediate dates are missing after frequency alignment (for example via \texttt{asfreq}), larger \texttt{period} comparisons can be skipped or shifted to different comparable dates.

\paragraph{Returns}
A \texttt{Series} indexed by date.

\subsubsection{\texttt{factor\_rank\_autocorrelation(factor\_data, period=1)}}
Measure the autocorrelation of factor ranks across periods.

\paragraph{Behavior}
\begin{itemize}
  \item Ranks assets by their factor values within each date, then computes the correlation between each date's rank vector and the rank vector period dates earlier.
  \item Useful as a turnover/stability diagnostic.
\end{itemize}

\paragraph{Worked example (\texttt{period=1})}
Suppose the factor values for assets \texttt{A}, \texttt{B}, \texttt{C}, \texttt{D} are:

\begin{center}
\begin{tabular}{lll}
date & factor values (A, B, C, D) & rank vector (A, B, C, D) \\
2026-01-02 & \((1, 2, 3, 4)\) & \((1, 2, 3, 4)\) \\
2026-01-03 & \((4, 3, 2, 1)\) & \((4, 3, 2, 1)\) \\
2026-01-06 & \((1, 2, 3, 4)\) & \((1, 2, 3, 4)\) \\
\end{tabular}
\end{center}

With \texttt{period=1}, each date is correlated with the immediately prior date:

\begin{itemize}
  \item On \texttt{2026-01-03}: corr\((4, 3, 2, 1), (1, 2, 3, 4)\) = \(-1.0\).
  \item On \texttt{2026-01-06}: corr\((1, 2, 3, 4), (4, 3, 2, 1)\) = \(-1.0\).
\end{itemize}

The first date has no prior rank vector, so its value is \texttt{NaN}.

\paragraph{Companion example (\texttt{period=2})}
Using the same data with \texttt{period=2}, \texttt{2026-01-06} is compared to \texttt{2026-01-02}:

\begin{itemize}
  \item corr\((1, 2, 3, 4), (1, 2, 3, 4)\) = \(1.0\).
\end{itemize}

So a larger \texttt{period} can reveal longer-horizon rank stability even when adjacent dates are unstable.

\paragraph{Returns}
A \texttt{Series} of autocorrelation values indexed by date.

\subsection{Event-study helpers}

\subsubsection{\texttt{common\_start\_returns(factor, returns, before, after, cumulative=False, mean\_by\_date=False, demean\_by=None)}}
Align return windows around common event dates.

\paragraph{Behavior}
\begin{itemize}
  \item Builds a return window around each factor date and asset pair.
  \item Aligns all windows to a common event-time index.
  \item Can work with cumulative or period returns.
  \item Can de-mean against a reference universe.
  \item Can average across assets by date.
\end{itemize}

\paragraph{Returns}
A \texttt{DataFrame} of aligned return windows.

\subsubsection{\texttt{average\_cumulative\_return\_by\_quantile(factor\_data, returns, periods\_before=10, periods\_after=15, demeaned=True, group\_adjust=False, by\_group=False)}}
Compute average cumulative returns around factor events by quantile.

\paragraph{Behavior}
\begin{itemize}
  \item Uses \texttt{common\_start\_returns()} internally.
  \item Computes mean and standard deviation of event-time cumulative returns.
  \item Can separate results by group.
  \item Supports group-neutral and universe-demeaned variants.
\end{itemize}

\paragraph{Returns}
A MultiIndex \texttt{DataFrame} containing mean and standard deviation across the event window.

\subsection{Portfolio simulation outputs}

\subsubsection{\texttt{factor\_cumulative\_returns(factor\_data, period, long\_short=True, group\_neutral=False, equal\_weight=False, quantiles=None, groups=None)}}
Simulate a factor portfolio and return cumulative performance.

\paragraph{Behavior}
\begin{itemize}
  \item Filters to a single forward-return horizon specified by \texttt{period}.
  \item Can limit analysis to selected quantiles or groups.
  \item Uses \texttt{factor\_returns()} and then converts the result to cumulative returns.
\end{itemize}

This means cumulative factor returns are built directly from the period-wise factor return series described above.

\paragraph{Returns}
A cumulative return \texttt{Series}.

\subsubsection{\texttt{factor\_positions(factor\_data, period, long\_short=True, group\_neutral=False, equal\_weight=False, quantiles=None, groups=None)}}
Simulate a factor portfolio and return the time series of positions.

\paragraph{Behavior}
\begin{itemize}
  \item Filters to the requested forward-return horizon.
  \item Reuses \texttt{factor\_weights()} to compute holdings.
  \item Converts weights into rolling positions with \texttt{positions()}.
\end{itemize}

\paragraph{Returns}
A \texttt{DataFrame} of asset positions over time.

\subsubsection{\texttt{create\_pyfolio\_input(factor\_data, period, capital=None, long\_short=True, group\_neutral=False, equal\_weight=False, quantiles=None, groups=None, benchmark\_period="1D")}}
Create returns, positions, and benchmark data in the format expected by Pyfolio.

\paragraph{Behavior}
\begin{itemize}
  \item Builds cumulative strategy returns.
  \item Resamples returns and positions to daily frequency.
  \item Adds a \texttt{cash} column to the positions output.
  \item Optionally converts percentage positions into dollar positions using \texttt{capital}.
  \item Computes a benchmark series from the factor universe when the benchmark period is available.
\end{itemize}

\paragraph{Returns}
A tuple: \((\text{returns}, \text{positions}, \text{benchmark})\).

\subsection{Dependencies and implementation notes}

\begin{itemize}
  \item Uses \texttt{pandas}, \texttt{numpy}, \texttt{scipy.stats}, \texttt{statsmodels}, and \texttt{empyrical}.
  \item Relies heavily on \texttt{alphalens.utils} for forward-return column detection and return demeaning.
  \item Most functions assume data has already been cleaned and aligned into the Alphalens MultiIndex format.
  \item Several functions use a trading-calendar frequency attached to the date index, so preserving index frequency is important for correct behavior.
\end{itemize}

\subsection{Practical usage}

Typical workflow:

\begin{enumerate}
  \item Build cleaned factor data with \texttt{alphalens.utils.get\_clean\_factor\_and\_forward\_returns()}.
  \item Compute diagnostics such as IC, mean returns by quantile, and turnover.
  \item Simulate factor-weighted returns or positions.
  \item Feed the outputs into plotting or Pyfolio workflows.
\end{enumerate}

Example call: \texttt{factor\_information\_coefficient(factor\_data)} or \texttt{create\_pyfolio\_input(factor\_data, period="1D")}.

\section{alphalens.plotting}

Visualization utilities for factor research, including plots for Information Coefficient analysis, factor returns, quantile performance, and other diagnostic charts. The module provides publication-quality visualizations with sane defaults for the Alphalens workflow.

This module visualizes the outputs from \texttt{alphalens.performance} and the cleaned factor data from \texttt{alphalens.utils}.

\subsection{Styling and context management}

\subsubsection{\texttt{customize(func)} --- decorator}
Apply default plotting context and style to a function.

\paragraph{Behavior}
\begin{itemize}
  \item Sets seaborn colorblind-friendly palette.
  \item Applies default plotting context and axes style.
  \item Can be disabled with \texttt{set\_context=False} in the decorated function call.
\end{itemize}

\subsubsection{\texttt{plotting\_context(context="notebook", font\_scale=1.5, rc=None)}}
Create an alphalens-specific plotting context.

\paragraph{Behavior}
\begin{itemize}
  \item Wraps \texttt{seaborn.plotting\_context()} with alphalens defaults.
  \item Default line width is 1.5.
  \item Can be used as a context manager for multiple plots.
\end{itemize}

\paragraph{Returns}
A seaborn plotting context object suitable for use in \texttt{with} statements.

\subsubsection{\texttt{axes\_style(style="darkgrid", rc=None)}}
Create an alphalens-specific axes style context.

\paragraph{Behavior}
\begin{itemize}
  \item Wraps \texttt{seaborn.axes\_style()} with alphalens defaults.
  \item Commonly used with \texttt{plotting\_context()}.
\end{itemize}

\paragraph{Returns}
A seaborn axes style object suitable for use in \texttt{with} statements.

\subsection{Data table displays}

\subsubsection{\texttt{plot\_returns\_table(alpha\_beta, mean\_ret\_quantile, mean\_ret\_spread\_quantile, return\_df=False)}}
Display a summary table of factor portfolio returns, alpha, beta, and quantile statistics.

\paragraph{Behavior}
\begin{itemize}
  \item Combines alpha/beta results with quantile return statistics.
  \item Converts decimal returns to basis points.
  \item Prints a formatted table or returns as a \texttt{DataFrame}.
\end{itemize}

\paragraph{Returns}
Pretty-printed table or a \texttt{DataFrame} when \texttt{return\_df=True}.

\subsubsection{\texttt{plot\_turnover\_table(autocorrelation\_data, quantile\_turnover, return\_df=False)}}
Display quantile turnover and factor rank autocorrelation statistics.

\paragraph{Behavior}
\begin{itemize}
  \item Shows mean turnover for each quantile by period.
  \item Shows mean factor rank autocorrelation.
  \item Prints formatted tables or returns as \texttt{DataFrame}s.
\end{itemize}

\paragraph{Returns}
Printed tables or a pair of \texttt{DataFrame}s when \texttt{return\_df=True}.

\subsubsection{\texttt{plot\_information\_table(ic\_data, return\_df=False)}}
Display Information Coefficient summary statistics.

\paragraph{Behavior}
\begin{itemize}
  \item Computes IC mean, standard deviation, risk-adjusted IC, and t-statistics.
  \item Includes skewness and kurtosis.
  \item Prints a summary table or returns as a \texttt{DataFrame}.
\end{itemize}

\paragraph{Returns}
Pretty-printed table or a \texttt{DataFrame} when \texttt{return\_df=True}.

\subsubsection{\texttt{plot\_quantile\_statistics\_table(factor\_data, return\_df=False)}}
Display factor value statistics by quantile.

\paragraph{Behavior}
\begin{itemize}
  \item Groups factor data by \texttt{factor\_quantile}.
  \item Reports min, max, mean, standard deviation, and count per quantile.
  \item Includes count percentages.
\end{itemize}

\paragraph{Returns}
Pretty-printed table or a \texttt{DataFrame} when \texttt{return\_df=True}.

\subsection{Information Coefficient visualization}

\subsubsection{\texttt{plot\_ic\_ts(ic, ax=None)}}
Plot the Information Coefficient time series with a rolling 22-day average.

\paragraph{Behavior}
\begin{itemize}
  \item Creates one subplot per forward-return period.
  \item Overlays daily IC and a 22-day moving average.
  \item Includes mean and standard deviation annotations.
  \item Aligns y-axis ranges across subplots.
\end{itemize}

\paragraph{Returns}
A numpy array of matplotlib axes.

\subsubsection{\texttt{plot\_ic\_hist(ic, ax=None)}}
Plot IC as a histogram with kernel density estimate.

\paragraph{Behavior}
\begin{itemize}
  \item Creates a 3-column grid of histograms, one per forward-return period.
  \item Includes KDE overlay and a vertical line at the mean.
  \item Clamps x-axis to \texttt{[-1, 1]}.
\end{itemize}

\paragraph{Returns}
A numpy array of matplotlib axes.

\subsubsection{\texttt{plot\_ic\_qq(ic, theoretical\_dist=scipy.stats.norm, ax=None)}}
Plot IC values against a theoretical distribution using a Q-Q plot.

\paragraph{Behavior}
\begin{itemize}
  \item Creates a 3-column grid of Q-Q plots.
  \item Supports normal and Student-t distributions.
  \item Uses \texttt{statsmodels.qqplot()} for plotting.
\end{itemize}

\paragraph{Returns}
A numpy array of matplotlib axes.

\subsubsection{\texttt{plot\_ic\_by\_group(ic\_group, ax=None)}}
Plot IC by group as a bar chart.

\paragraph{Behavior}
\begin{itemize}
  \item Shows mean IC for each group.
  \item Rotates x-axis labels for readability.
\end{itemize}

\paragraph{Returns}
A matplotlib axes object.

\subsubsection{\texttt{plot\_monthly\_ic\_heatmap(mean\_monthly\_ic, ax=None)}}
Plot IC or returns as a month/year heatmap with color coding.

\paragraph{Behavior}
\begin{itemize}
  \item Creates one heatmap per forward-return period.
  \item Uses reversed coolwarm colormap (red = negative, blue = positive).
  \item Annotates cell values.
  \item Centered at 0.
\end{itemize}

\paragraph{Returns}
A numpy array of matplotlib axes.

\subsection{Returns analysis}

\subsubsection{\texttt{plot\_quantile\_returns\_bar(mean\_ret\_by\_q, by\_group=False, ylim\_percentiles=None, ax=None)}}
Plot mean returns by factor quantile as a bar chart.

\paragraph{Behavior}
\begin{itemize}
  \item Optionally creates separate subplots for each group.
  \item Can limit y-axis range using percentiles of observed data.
  \item Converts to basis points.
\end{itemize}

\paragraph{Returns}
A matplotlib axes object or array depending on grouping.

\subsubsection{\texttt{plot\_quantile\_returns\_violin(return\_by\_q, ylim\_percentiles=None, ax=None)}}
Plot return distributions by factor quantile using violin plots.

\paragraph{Behavior}
\begin{itemize}
  \item Shows distribution of returns within each quantile across periods.
  \item Uses seaborn violin plot with quartile indicators.
  \item Optionally limits y-axis range using percentiles.
\end{itemize}

\paragraph{Returns}
A matplotlib axes object.

\subsubsection{\texttt{plot\_mean\_quantile\_returns\_spread\_time\_series(mean\_returns\_spread, std\_err=None, bandwidth=1, ax=None)}}
Plot the top-minus-bottom quantile return spread over time.

\paragraph{Behavior}
\begin{itemize}
  \item Shows the difference between top and bottom quantile returns period by period.
  \item Overlays a 22-day rolling average.
  \item Optionally adds confidence bands based on standard error.
  \item Symmetrizes y-axis around zero at the 95th percentile.
\end{itemize}

\paragraph{Returns}
A matplotlib axes object or array if input is a \texttt{DataFrame}.

\subsubsection{\texttt{plot\_cumulative\_returns(factor\_returns, period, freq=None, title=None, ax=None)}}
Plot cumulative returns of a factor-weighted portfolio.

\paragraph{Behavior}
\begin{itemize}
  \item Converts simple returns to cumulative using \texttt{performance.cumulative\_returns()}.
  \item Includes a reference line at 1.0 (break-even).
  \item Displays the forward-return period in the title.
\end{itemize}

\paragraph{Returns}
A matplotlib axes object.

\subsubsection{\texttt{plot\_cumulative\_returns\_by\_quantile(quantile\_returns, period, freq=None, ax=None)}}
Plot cumulative returns separately for each factor quantile.

\paragraph{Behavior}
\begin{itemize}
  \item Uses reversed coolwarm colormap for quantile coloring.
  \item Uses symmetric log scale on the y-axis for readability across wide ranges.
  \item Shows all quantiles on a single plot with a legend.
\end{itemize}

\paragraph{Returns}
A matplotlib axes object.

\subsubsection{\texttt{plot\_quantile\_average\_cumulative\_return(avg\_cumulative\_returns, by\_quantile=False, std\_bar=False, title=None, ax=None)}}
Plot average cumulative returns by factor quantile around event dates.

\paragraph{Behavior}
\begin{itemize}
  \item Shows mean cumulative returns across the event window for each quantile.
  \item Optionally creates separate subplots per quantile for clarity.
  \item Can overlay standard deviation error bars.
  \item Uses reversed coolwarm colormap for consistency.
  \item Marks event date (x=0) with a vertical dashed line.
\end{itemize}

\paragraph{Returns}
A matplotlib axes object or array.

\subsection{Factor characteristics}

\subsubsection{\texttt{plot\_factor\_rank\_auto\_correlation(factor\_autocorrelation, period=1, ax=None)}}
Plot factor rank autocorrelation over time.

\paragraph{Behavior}
\begin{itemize}
  \item Shows the stability of relative factor rankings.
  \item Includes a reference line at 0 (no autocorrelation).
  \item Annotates the mean autocorrelation in a text box.
\end{itemize}

\paragraph{Returns}
A matplotlib axes object.

\subsubsection{\texttt{plot\_top\_bottom\_quantile\_turnover(quantile\_turnover, period=1, ax=None)}}
Plot turnover for the top and bottom quantiles over time.

\paragraph{Behavior}
\begin{itemize}
  \item Extracts the maximum and minimum quantiles.
  \item Shows the proportion of names new to each quantile each period.
  \item Useful for detecting data quality or stationarity issues.
\end{itemize}

\paragraph{Returns}
A matplotlib axes object.

\subsubsection{\texttt{plot\_events\_distribution(events, num\_bars=50, ax=None)}}
Plot the distribution of factor events across time.

\paragraph{Behavior}
\begin{itemize}
  \item Divides the time span into \texttt{num\_bars} intervals.
  \item Counts events in each interval.
  \item Displays as a bar chart.
\end{itemize}

\paragraph{Returns}
A matplotlib axes object.

\subsection{Constants}

\subsubsection{\texttt{DECIMAL\_TO\_BPS = 10000}}
Conversion factor for decimal returns to basis points. Used throughout the module when displaying returns-related plots.

\subsection{Dependencies and conventions}

\begin{itemize}
  \item \textbf{Visualization}: seaborn, matplotlib, scipy.stats
  \item \textbf{Analysis}: pandas, numpy, statsmodels
  \item \textbf{Internal}: \texttt{alphalens.utils}, \texttt{alphalens.performance}
\end{itemize}

\paragraph{Color conventions}
\begin{itemize}
  \item \textbf{Coolwarm reversed}: Red represents low/negative quantiles; blue represents high/positive quantiles.
  \item \textbf{Colorblind palette}: All plots use seaborn's colorblind palette for accessibility.
  \item \textbf{Line colors}: Green (forestgreen) for means and averages; blue (steelblue) for confidence bands; red (orangered) for rolling averages.
\end{itemize}

\paragraph{Plot conventions}
\begin{itemize}
  \item \textbf{Multi-subplot layouts}: Plots with multiple periods or groups create grid layouts (usually 1\(\times\)3 or 2\(\times\)2).
  \item \textbf{Shared axes}: Where applicable, y-axes are shared across subplots to aid comparison.
  \item \textbf{Reference lines}: Zero lines are included on relevant plots to indicate break-even or neutral positions.
  \item \textbf{Log scales}: Cumulative returns plots use symmetric log scale (\texttt{symlog}) to handle very large or very small returns.
  \item \textbf{Annotations}: Statistical summaries (mean, std dev) are included in text boxes on many plots.
\end{itemize}

\subsection{Practical usage}

Typical workflow:

\begin{enumerate}
  \item Compute factor diagnostics with \texttt{alphalens.performance} functions.
  \item Call plotting functions with the results.
  \item All plotting functions are decorated with \texttt{@customize} to apply consistent styling.
  \item Optional: use \texttt{set\_context=False} in any plotting call to skip automatic styling.
\end{enumerate}

Example call: \texttt{plot\_ic\_ts(ic\_data)} or \texttt{plot\_quantile\_returns\_bar(mean\_ret\_by\_q)}.

\subsection{Notes}

\begin{itemize}
  \item All plotting functions accept an optional \texttt{ax} parameter for custom subplot layouts.
  \item Return values vary: some functions print tables directly, others return \texttt{DataFrame}s or axes.
  \item The \texttt{return\_df=False} parameter on table functions controls printing vs. returning.
  \item Most functions are designed for Jupyter notebooks but work in any matplotlib environment.
\end{itemize}

\section{alphalens.tears}

Tear sheet generation for comprehensive factor analysis. A tear sheet is a multi-plot report that visualizes factor performance from multiple perspectives, including returns analysis, information coefficient statistics, turnover metrics, and event studies.

This module orchestrates the workflows of \texttt{alphalens.performance} and \texttt{alphalens.plotting} to produce production-quality analysis reports suitable for presentation and decision-making.

\subsection{Helper class}

\subsubsection{\texttt{GridFigure}}
Utility class for managing matplotlib grid-based subplot layouts with automatic row/column navigation.

\paragraph{Methods}
\begin{itemize}
  \item \texttt{\_\_init\_\_(rows, cols)} --- Create a grid figure with specified dimensions.
  \item \texttt{next\_row()} --- Get the next full-width subplot (advances to the next row).
  \item \texttt{next\_cell()} --- Get the next cell in the current row (auto-wraps to next row when full).
  \item \texttt{close()} --- Close the figure and clean up resources.
\end{itemize}

\paragraph{Behavior}
\begin{itemize}
  \item Default figure size is 14 inches wide by 7 inches per row.
  \item Horizontal and vertical spacing are set to 0.4 and 0.3 respectively for readability.
\end{itemize}

\subsection{Tear sheet generators}

All tear sheet functions are decorated with \texttt{@plotting.customize} to apply consistent styling automatically. The \texttt{set\_context=False} parameter can be passed to any tear sheet to skip the decorator.

\subsubsection{\texttt{create\_summary\_tear\_sheet(factor\_data, long\_short=True, group\_neutral=False)}}
Generate a compact tear sheet with key returns, information, and turnover statistics.

\paragraph{Behavior}
\begin{itemize}
  \item Displays factor quantile statistics table.
  \item Computes and shows alpha, beta, top/bottom quantile returns, and spreads.
  \item Plots factor quantile returns bar chart.
  \item Computes and displays IC summary statistics.
  \item Analyzes turnover and factor rank autocorrelation.
  \item All metrics respect the \texttt{long\_short} and \texttt{group\_neutral} flags.
\end{itemize}

\paragraph{Returns}
Displays a matplotlib figure with the complete summary.

\subsubsection{\texttt{create\_returns\_tear\_sheet(factor\_data, long\_short=True, group\_neutral=False, by\_group=False)}}
Detailed tear sheet focused on returns analysis by factor quantile.

\paragraph{Behavior}
\begin{itemize}
  \item Computes factor portfolio returns.
  \item Displays alpha and beta metrics.
  \item Shows mean quantile returns as a bar chart.
  \item Plots return distributions using violin plots (by date across periods).
  \item If \texttt{'1D'} returns are available, plots cumulative portfolio returns and quantile-level cumulative returns.
  \item Displays top-minus-bottom quantile return spread over time with optional error bands.
  \item When \texttt{by\_group=True}, creates separate bar charts for each asset group.
\end{itemize}

\paragraph{Parameters}
\begin{itemize}
  \item \texttt{long\_short} --- Enable long-short demeaning for dollar-neutral analysis.
  \item \texttt{group\_neutral} --- Normalize weights and returns across groups.
  \item \texttt{by\_group} --- Create separate visualizations per group.
\end{itemize}

\paragraph{Returns}
Displays one or more matplotlib figures.

\subsubsection{\texttt{create\_information\_tear\_sheet(factor\_data, group\_neutral=False, by\_group=False)}}
Tear sheet focused on Information Coefficient (IC) analysis.

\paragraph{Behavior}
\begin{itemize}
  \item Displays IC summary statistics (mean, std, risk-adjusted IC, t-stat, p-value, skewness, kurtosis).
  \item Plots IC time series with 22-day rolling average for each forward-return period.
  \item Plots IC histograms with KDE.
  \item Plots IC Q-Q plots against normal distribution.
  \item When \texttt{by\_group=False} (default), creates a heatmap of monthly mean IC.
  \item When \texttt{by\_group=True}, shows mean IC separately for each group.
\end{itemize}

\paragraph{Parameters}
\begin{itemize}
  \item \texttt{group\_neutral} --- Demean forward returns by group before computing IC.
  \item \texttt{by\_group} --- Separate analysis by asset group.
\end{itemize}

\paragraph{Returns}
Displays matplotlib figure(s).

\subsubsection{\texttt{create\_turnover\_tear\_sheet(factor\_data, turnover\_periods=None)}}
Tear sheet analyzing factor turnover and rank stability.

\paragraph{Behavior}
\begin{itemize}
  \item Displays quantile turnover statistics (mean turnover per quantile by period).
  \item Displays factor rank autocorrelation statistics.
  \item Plots time series of top and bottom quantile turnover for each period.
  \item Plots factor rank autocorrelation over time for each period.
  \item Automatically extracts day-multiple forward-return periods if \texttt{turnover\_periods} is not provided.
\end{itemize}

\paragraph{Parameters}
\begin{itemize}
  \item \texttt{turnover\_periods} --- Custom periods for turnover analysis. If not provided, uses exact day multiples from \texttt{factor\_data}.
\end{itemize}

\paragraph{Returns}
Displays matplotlib figure with turnover and autocorrelation plots.

\subsubsection{\texttt{create\_full\_tear\_sheet(factor\_data, long\_short=True, group\_neutral=False, by\_group=False)}}
Comprehensive tear sheet combining all analysis perspectives.

\paragraph{Behavior}
\begin{itemize}
  \item Displays factor quantile statistics.
  \item Calls \texttt{create\_returns\_tear\_sheet()} with all parameters passed through.
  \item Calls \texttt{create\_information\_tear\_sheet()} with appropriate flags.
  \item Calls \texttt{create\_turnover\_tear\_sheet()} with defaults.
  \item All component tear sheets skip automatic styling via \texttt{set\_context=False} to apply consistent styling once.
\end{itemize}

\paragraph{Parameters}
\begin{itemize}
  \item \texttt{long\_short} --- Enable long-short demeaning.
  \item \texttt{group\_neutral} --- Enable group-neutral normalization.
  \item \texttt{by\_group} --- Separate analysis per group.
\end{itemize}

\paragraph{Returns}
Displays multiple matplotlib figures for all analyses.

\subsubsection{\texttt{create\_event\_returns\_tear\_sheet(factor\_data, returns, avgretplot=(5, 15), long\_short=True, group\_neutral=False, std\_bar=True, by\_group=False)}}
Tear sheet analyzing average cumulative returns around factor events (event study).

\paragraph{Behavior}
\begin{itemize}
  \item Computes average cumulative returns within a window before and after each factor event.
  \item Plots overall average cumulative returns with optional error bands.
  \item When \texttt{std\_bar=True}, plots separate cumulative return traces for each quantile with standard deviation error bars.
  \item When \texttt{by\_group=True}, repeats the analysis separately for each asset group.
  \item Useful for analyzing the predictive power of a factor over different time horizons relative to the event date.
\end{itemize}

\paragraph{Parameters}
\begin{itemize}
  \item \texttt{factor\_data} --- MultiIndex factor data (no forward returns required for this tear sheet).
  \item \texttt{returns} --- DataFrame of asset returns indexed by date.
  \item \texttt{avgretplot} --- Tuple of \texttt{(periods\_before, periods\_after)} for the event window.
  \item \texttt{long\_short} --- Demean cumulative returns.
  \item \texttt{group\_neutral} --- Normalize returns by group.
  \item \texttt{std\_bar} --- Include error bar plots by quantile.
  \item \texttt{by\_group} --- Separate analysis per group.
\end{itemize}

\paragraph{Returns}
Displays matplotlib figure(s) showing event-window average returns.

\subsubsection{\texttt{create\_event\_study\_tear\_sheet(factor\_data, returns, avgretplot=(5, 15), rate\_of\_ret=True, n\_bars=50)}}
Tear sheet for analyzing specific events or signals.

\paragraph{Behavior}
\begin{itemize}
  \item Displays factor quantile statistics.
  \item Plots histogram of event distribution across time (divided into \texttt{n\_bars} intervals).
  \item Calls \texttt{create\_event\_returns\_tear\_sheet()} internally for average cumulative returns analysis.
  \item Computes and displays mean quantile returns.
  \item Plots return distributions via violin plots by quantile.
  \item Can optionally annualize returns using \texttt{rate\_of\_ret=True}.
\end{itemize}

\paragraph{Parameters}
\begin{itemize}
  \item \texttt{factor\_data} --- MultiIndex factor data indexed by event dates.
  \item \texttt{returns} --- Asset returns DataFrame.
  \item \texttt{avgretplot} --- Window for event-study cumulative returns.
  \item \texttt{rate\_of\_ret} --- Convert simple returns to annualized rates.
  \item \texttt{n\_bars} --- Number of time intervals for event distribution histogram.
\end{itemize}

\paragraph{Returns}
Displays matplotlib figures for event analysis.

\subsection{Data model}

All tear sheet functions expect \texttt{factor\_data} to be a MultiIndex \texttt{DataFrame} with:

\begin{itemize}
  \item Index: \texttt{(date, asset)}
  \item Columns:
  \begin{itemize}
    \item Forward-return columns (e.g., \texttt{1D}, \texttt{5D}, \texttt{10D})
    \item \texttt{factor} --- The factor values
    \item \texttt{factor\_quantile} --- Quantile assignments
    \item \texttt{group} (optional) --- Asset group identifiers
  \end{itemize}
\end{itemize}

This is the standard output format from \texttt{alphalens.utils.get\_clean\_factor\_and\_forward\_returns()}.

\subsection{Dependencies and concepts}

\begin{itemize}
  \item \textbf{Visualization}: matplotlib, seaborn
  \item \textbf{Analysis}: pandas, numpy, scipy.stats, statsmodels
  \item \textbf{Internal}: \texttt{alphalens.plotting}, \texttt{alphalens.performance}, \texttt{alphalens.utils}
\end{itemize}

\paragraph{Key tear sheet concepts}
\begin{itemize}
  \item \textbf{Long-short analysis}: Factor values are demeaned, creating equal-weight long and short sides.
  \item \textbf{Group-neutral analysis}: Asset groups are weighted equally, and returns are normalized within groups.
  \item \textbf{Quantile analysis}: Assets are divided into quantiles/buckets based on factor values; returns and metrics are computed per quantile.
  \item \textbf{Information Coefficient}: Spearman rank correlation between factor values and forward returns; higher IC indicates stronger predictive power.
  \item \textbf{Turnover}: Proportion of assets entering/leaving a quantile each period; low turnover suggests stable rankings.
  \item \textbf{Event study}: Analysis of returns in a window around specific factor events or signals.
\end{itemize}

\subsection{Practical usage}

Typical workflow:

\begin{enumerate}
  \item Prepare cleaned factor data with \texttt{alphalens.utils.get\_clean\_factor\_and\_forward\_returns()}.
  \item Call the desired tear sheet function.
  \item For a quick overview, use \texttt{create\_summary\_tear\_sheet()}.
  \item For detailed analysis, use \texttt{create\_full\_tear\_sheet()}.
  \item For event-based analysis, use \texttt{create\_event\_returns\_tear\_sheet()} or \texttt{create\_event\_study\_tear\_sheet()}.
\end{enumerate}

Example:

\begin{verbatim}
from alphalens.tears import create_full_tear_sheet

create_full_tear_sheet(factor_data, long_short=True, group_neutral=False)
\end{verbatim}

\subsection{Notes}

\begin{itemize}
  \item All tear sheets automatically apply consistent styling via the \texttt{@plotting.customize} decorator.
  \item GridFigure handles the layout complexity automatically; most workflows don't require direct use.
  \item The \texttt{by\_group} parameter enables side-by-side or panel analysis when asset groups are present.
  \item For large datasets or many forward-return periods, tear sheets can take time to render; \texttt{create\_summary\_tear\_sheet()} is faster for quick feedback.
  \item All plotting occurs within the tear sheet functions; call \texttt{plt.show()} or allow the Jupyter kernel to display results automatically.
\end{itemize}

\end{document}

